\shead{Thể tích và đơn vị đo thể tích}
\gdef\cuboidVT#1#2#3#4{\coordinate (Q) at (0,0); \coordinate (P) at (#1,0); \coordinate (C) at (#1,#3); \coordinate (D) at (0,#3);
 \coordinate (M) at ({0.5*#2},{0.38*#2}); \coordinate (N) at ($(P)+(M)$); \coordinate (B) at ($(C)+(M)$); \coordinate (A) at ($(D)+(M)$);
 \fill[tmBlue!6] (Q)--(P)--(C)--(D)--cycle; \fill[tmBlue!3] (D)--(C)--(B)--(A)--cycle; \fill[tmBlue!13] (P)--(N)--(B)--(C)--cycle;
 \ifdim #4pt>0pt \coordinate (wl) at (0,#4);
  \fill[cyan!30] (Q)--(P)--($(P)+(wl)$)--($(Q)+(wl)$)--cycle; \fill[cyan!45] (P)--(N)--($(N)+(wl)$)--($(P)+(wl)$)--cycle;
  \fill[cyan!16] ($(Q)+(wl)$)--($(P)+(wl)$)--($(N)+(wl)$)--($(M)+(wl)$)--cycle;
  \draw[cyan!60!black,line width=0.5pt] ($(Q)+(wl)$)--($(P)+(wl)$)--($(N)+(wl)$);
  \draw[cyan!60!black,line width=0.4pt,dashed] ($(N)+(wl)$)--($(M)+(wl)$)--($(Q)+(wl)$); \fi
 \draw[geo,dashed,line width=0.6pt] (M)--(Q) (M)--(N) (M)--(A);
 \draw[geo] (Q)--(P)--(C)--(D)--cycle (D)--(A)--(B)--(C) (P)--(N)--(B);}
\begin{knowbox}
\begin{minipage}[t]{0.6\linewidth}
\textbf{Thể tích} của một hình là độ lớn phần không gian mà hình đó chiếm chỗ. Để đo thể tích, ta đếm xem hình đó chứa bao nhiêu \emph{hình lập phương đơn vị}.\par
\smallskip
\textbf{Xăng-ti-mét khối} (cm$^3$) là thể tích của hình lập phương có cạnh dài $1$ cm. Tương tự, $1$ dm$^3$ là thể tích hình lập phương cạnh $1$ dm, $1$ m$^3$ là thể tích hình lập phương cạnh $1$ m.
\[1\text{ m}^3=1\,000\text{ dm}^3\qquad 1\text{ dm}^3=1\,000\text{ cm}^3\]
Mỗi đơn vị đo thể tích gấp $1\,000$ lần đơn vị bé hơn tiếp liền, nên khi đổi đơn vị ta dời dấu phẩy \textbf{ba} chữ số.\par
\smallskip
Với chất lỏng: \quad $1$~l $=1$ dm$^3$, \quad $1$~ml $=1$ cm$^3$, \quad \mbox{$1$ m$^3=1\,000$~l}.
\end{minipage}\hfill
\begin{minipage}[t]{0.36\linewidth}\vspace{6pt}\centering
\begin{tikzpicture}[scale=0.62]
\cuboidVT{3}{3}{3}{0}
\foreach \i in {1,2}{\draw[tmBlue!60,line width=0.35pt] (\i,0)--(\i,3)--++(1.5,1.14) (0,\i)--(3,\i)--++(1.5,1.14)
 ($(0,3)+\i*(0.5,0.38)$)--++(3,0)--++(0,-3);}
\draw[geo] (Q)--(P)--(C)--(D)--cycle (D)--(A)--(B)--(C) (P)--(N)--(B);
\node[dim,below] at ($(Q)!0.5!(P)$){$1$ dm $=10$ cm};
\end{tikzpicture}\par\smallskip
{\scriptsize\color{Slate} $1$ dm$^3$: $10$ lớp, mỗi lớp $10\times10$ hình lập phương $1$ cm$^3$, tức là $1\,000$ cm$^3$ (hình vẽ thu gọn)}
\end{minipage}
\end{knowbox}
\shead{Thể tích hình hộp chữ nhật, hình lập phương}
\begin{knowbox}
\begin{minipage}[t]{0.6\linewidth}
Hình hộp chữ nhật dài $5$ cm, rộng $3$ cm, cao $2$ cm xếp được $2$ lớp, mỗi lớp $5\times3=15$ hình lập phương $1$ cm$^3$, tức là $5\times3\times2=30$ cm$^3$.\par
\smallskip
Muốn tính \textbf{thể tích hình hộp chữ nhật}, ta lấy chiều dài nhân với chiều rộng rồi nhân với chiều cao (cùng một đơn vị đo):
\[V=a\times b\times c=S_{\text{đáy}}\times c\]
\textbf{Thể tích hình lập phương} cạnh $a$: \quad $V=a\times a\times a$.\par
\smallskip
Tính ngược lại: \quad $c=V:(a\times b)$, \quad $S_{\text{đáy}}=V:c$.
\end{minipage}\hfill
\begin{minipage}[t]{0.36\linewidth}\vspace{6pt}\centering
\begin{tikzpicture}[scale=0.55]
\cuboidVT{5}{3}{2}{0}
\foreach \i in {1,...,4}{\draw[tmBlue!60,line width=0.35pt] (\i,0)--(\i,2)--++(1.5,1.14);}
\draw[tmBlue!60,line width=0.35pt] (0,1)--(5,1)--++(1.5,1.14);
\foreach \k in {1,2}{\draw[tmBlue!60,line width=0.35pt] ($(0,2)+\k*(0.5,0.38)$)--++(5,0)--++(0,-2);}
\draw[geo] (Q)--(P)--(C)--(D)--cycle (D)--(A)--(B)--(C) (P)--(N)--(B);
\node[dim,below] at ($(Q)!0.5!(P)$){$5$ cm}; \node[dim,right] at ($(P)!0.5!(N)$){$3$ cm}; \node[dim,left] at ($(Q)!0.5!(D)$){$2$ cm};
\end{tikzpicture}
\end{minipage}
\end{knowbox}
\shead{Thể tích vật chìm trong nước}
\begin{knowbox}
Hòn đá, cây san hô có hình dạng không xác định nên không tính thể tích bằng công thức được. Thả vật vào bể nước sao cho vật \textbf{ngập hoàn toàn} trong nước: mực nước dâng lên vì vật đã chiếm một phần thể tích của bể.
\begin{itemize}[leftmargin=1.2em,itemsep=1pt,topsep=2pt]
\item \emph{Cách 1:} thể tích vật $=$ thể tích (nước $+$ vật) $-$ thể tích nước.
\item \emph{Cách 2:} thể tích vật $=$ diện tích đáy bể $\times$ chiều cao phần nước dâng lên.
\end{itemize}
Nếu bể đã đầy nước thì thể tích nước \textbf{tràn ra} bằng thể tích của vật.
\end{knowbox}
\gap
\begin{exbox}{\ExWord\ 1}
Viết số thích hợp vào chỗ chấm.
\begin{cols2}
\pa{a} $3{,}5$ m$^3=\ldots$ dm$^3$ & \pa{b} $2\,400$ cm$^3=\ldots$ dm$^3$\\
\pa{c} $0{,}75$ dm$^3=\ldots$ cm$^3=\ldots$ ml & \pa{d} $1{,}2$ m$^3=\ldots$ l
\end{cols2}
\guide
\pa{a} $1$ m$^3=1\,000$ dm$^3$ nên $3{,}5$ m$^3=3{,}5\times1\,000=3\,500$ dm$^3$.\par
\pa{b} $2\,400$ cm$^3=2\,400:1\,000=2{,}4$ dm$^3$.\par
\pa{c} $0{,}75$ dm$^3=750$ cm$^3=750$ ml.\par
\pa{d} $1{,}2$ m$^3=1\,200$ dm$^3=1\,200$ l.
\dapso{\pa{a} $3\,500$; \pa{b} $2{,}4$; \pa{c} $750$; $750$; \pa{d} $1\,200$.}
\end{exbox}
\gap
\begin{exbox}{\ExWord\ 2}
Tính thể tích của:\par
\pa{a} hình hộp chữ nhật có chiều dài $1{,}2$ m, chiều rộng $5$ dm, chiều cao $8$ dm;\par
\pa{b} hình lập phương có cạnh $2{,}5$ cm.
\guide
\pa{a} Đổi về cùng đơn vị đo: $1{,}2$ m $=12$ dm.\par $V=12\times5\times8=480$ dm$^3$.\par
\pa{b} $V=2{,}5\times2{,}5\times2{,}5=15{,}625$ cm$^3$.
\dapso{\pa{a} $480$ dm$^3$; \pa{b} $15{,}625$ cm$^3$.}
\end{exbox}
\gap
\begin{exbox}{\ExWord\ 3}
Một bể nước dạng hình hộp chữ nhật dài $1{,}5$ m, rộng $0{,}8$ m, cao $1$ m.\par
\pa{a} Khi đầy, bể chứa được bao nhiêu lít nước?\par
\pa{b} Hiện nay lượng nước trong bể bằng $\dfrac34$ thể tích bể. Trong bể có bao nhiêu lít nước?
\guide
\pa{a} $V=1{,}5\times0{,}8\times1=1{,}2$ m$^3=1\,200$ dm$^3=1\,200$ l.\par
\pa{b} $1\,200\times\dfrac34=900$ l.
\dapso{\pa{a} $1\,200$ l; \pa{b} $900$ l.}
\end{exbox}
\gap
\begin{exbox}{\ExWord\ 4}
Một bể kính có đáy dài $25$ cm, rộng $16$ cm. Mực nước trong bể cao $10$ cm. Sau khi thả vào bể một hòn đá (ngập hoàn toàn trong nước), mực nước cao $12{,}5$ cm. Tính thể tích hòn đá.
\begin{center}
\begin{tikzpicture}[scale=0.12]
\cuboidVT{25}{16}{16}{10}
\node[dim,below] at ($(Q)!0.5!(P)$){$25$ cm}; \node[dim,right] at ($(P)!0.5!(N)$){$16$ cm}; \node[dim,left] at (0,5){$10$ cm};
\begin{scope}[shift={(52,0)}]
\cuboidVT{25}{16}{16}{12.5}
\fill[Slate!75] (13,4.5) ellipse (4.5 and 2.6);
\node[dim,below] at ($(Q)!0.5!(P)$){$25$ cm}; \node[dim,right] at ($(P)!0.5!(N)$){$16$ cm}; \node[dim,left] at (0,6.25){$12{,}5$ cm};
\end{scope}
\end{tikzpicture}
\end{center}
\guide
\emph{Cách 1.} Thể tích nước: $25\times16\times10=4\,000$ cm$^3$.\par
Thể tích nước và hòn đá: $25\times16\times12{,}5=5\,000$ cm$^3$.\par
Thể tích hòn đá: $5\,000-4\,000=1\,000$ cm$^3$.\par
\emph{Cách 2.} Chiều cao phần nước dâng lên: $12{,}5-10=2{,}5$ cm.\par
Thể tích hòn đá: $25\times16\times2{,}5=1\,000$ cm$^3$.
\dapso{$1\,000$ cm$^3$ ($=1$ dm$^3$).}
\end{exbox}
\gap
\begin{exbox}{\ExWord\ 5}
Một hình lập phương có diện tích toàn phần $864$ cm$^2$. Tính thể tích hình lập phương đó.
\guide
Diện tích một mặt: $864:6=144$ cm$^2$. Vì $144=12\times12$ nên cạnh dài $12$ cm.\par
$V=12\times12\times12=1\,728$ cm$^3$.
\dapso{$1\,728$ cm$^3$.}
\end{exbox}
\dividerbar{Luyện tập}
\begin{probbox}{1}{\relax}
Viết số thích hợp vào chỗ chấm.
\begin{cols3}
\pa{a} $5$ m$^3=\ldots$ dm$^3$ & \pa{b} $3{,}2$ dm$^3=\ldots$ cm$^3$ & \pa{c} $4\,500$ cm$^3=\ldots$ dm$^3$\\
\pa{d} $750$ dm$^3=\ldots$ m$^3$ & \pa{e} $2{,}05$ m$^3=\ldots$ dm$^3$ & \pa{f} $\dfrac14$ m$^3=\ldots$ dm$^3$
\end{cols3}
\end{probbox}
\probsol{
\pa{a} $5\times1\,000=5\,000$.\quad \pa{b} $3{,}2\times1\,000=3\,200$.\quad \pa{c} $4\,500:1\,000=4{,}5$.\par
\pa{d} $750:1\,000=0{,}75$.\quad \pa{e} $2{,}05\times1\,000=2\,050$.\quad \pa{f} $1\,000:4=250$.
\dapso{$5\,000$;\ $3\,200$;\ $4{,}5$;\ $0{,}75$;\ $2\,050$;\ $250$.}}
\gap
\begin{probbox}{2}{\relax}
Viết số thích hợp vào chỗ chấm.
\begin{cols2}
\pa{a} $3{,}6$ m$^3=\ldots$ l & \pa{b} $450$ l $=\ldots$ m$^3$\\
\pa{c} $2{,}5$ l $=\ldots$ cm$^3$ & \pa{d} $750$ ml $=\ldots$ dm$^3$
\end{cols2}
\end{probbox}
\probsol{
\pa{a} $3{,}6$ m$^3=3\,600$ dm$^3=3\,600$ l.\par
\pa{b} $450$ l $=450$ dm$^3=0{,}45$ m$^3$.\par
\pa{c} $2{,}5$ l $=2{,}5$ dm$^3=2\,500$ cm$^3$.\par
\pa{d} $750$ ml $=750$ cm$^3=0{,}75$ dm$^3$.
\dapso{$3\,600$;\ $0{,}45$;\ $2\,500$;\ $0{,}75$.}}
\gap
\begin{probbox}{3}{\relax}
Điền dấu $>$, $<$, $=$ thích hợp vào chỗ chấm.
\begin{cols3}
\pa{a} $2{,}5$ m$^3\ \ldots\ 2\,500$ dm$^3$ & \pa{b} $3$ dm$^3$ $40$ cm$^3\ \ldots\ 3{,}4$ dm$^3$ & \pa{c} $0{,}9$ m$^3\ \ldots\ 90\,000$ cm$^3$
\end{cols3}
\end{probbox}
\probsol{
\pa{a} $2{,}5$ m$^3=2\,500$ dm$^3$ nên hai số đo bằng nhau.\par
\pa{b} $3$ dm$^3$ $40$ cm$^3=3\,040$ cm$^3$ và $3{,}4$ dm$^3=3\,400$ cm$^3$; $3\,040<3\,400$.\par
\pa{c} $0{,}9$ m$^3=900$ dm$^3=900\,000$ cm$^3>90\,000$ cm$^3$.
\dapso{\pa{a} $=$; \pa{b} $<$; \pa{c} $>$.}}
\gap
\begin{probbox}{4}{\relax}
Tính thể tích hình hộp chữ nhật có chiều dài $a$, chiều rộng $b$, chiều cao $c$:
\begin{cols2}
\pa{a} $a=8$ cm, $b=5$ cm, $c=4$ cm; & \pa{b} $a=1{,}5$ m, $b=0{,}8$ m, $c=1{,}2$ m;\\
\pa{c} $a=\dfrac23$ m, $b=\dfrac35$ m, $c=\dfrac12$ m; & \pa{d} $a=2$ dm, $b=15$ cm, $c=8$ cm.
\end{cols2}
\end{probbox}
\probsol{
\pa{a} $8\times5\times4=160$ cm$^3$.\par
\pa{b} $1{,}5\times0{,}8\times1{,}2=1{,}44$ m$^3$.\par
\pa{c} $\dfrac23\times\dfrac35\times\dfrac12=\dfrac15$ m$^3$.\par
\pa{d} $2$ dm $=20$ cm; $20\times15\times8=2\,400$ cm$^3$.
\dapso{$160$ cm$^3$; $1{,}44$ m$^3$; $\dfrac15$ m$^3$; $2\,400$ cm$^3$.}}
\gap
\begin{probbox}{5}{\relax}
Tính thể tích hình lập phương có cạnh:
\begin{cols3}
\pa{a} $7$ cm; & \pa{b} $1{,}2$ m; & \pa{c} $\dfrac34$ dm.
\end{cols3}
\end{probbox}
\probsol{
\pa{a} $7\times7\times7=343$ cm$^3$.\par
\pa{b} $1{,}2\times1{,}2\times1{,}2=1{,}728$ m$^3$.\par
\pa{c} $\dfrac34\times\dfrac34\times\dfrac34=\dfrac{27}{64}$ dm$^3$.
\dapso{$343$ cm$^3$; $1{,}728$ m$^3$; $\dfrac{27}{64}$ dm$^3$.}}
\gap
\begin{probbox}{6}{\relax}
\begin{minipage}[t]{0.58\linewidth}
Hình hộp bên được xếp kín bởi các hình lập phương cạnh $1$ cm. Có bao nhiêu hình lập phương nhỏ? Thể tích của hình hộp là bao nhiêu?
\end{minipage}\hfill
\begin{minipage}[t]{0.38\linewidth}\vspace{0pt}\centering
\begin{tikzpicture}[scale=0.42]
\cuboidVT{6}{4}{3}{0}
\foreach \i in {1,...,5}{\draw[tmBlue!60,line width=0.35pt] (\i,0)--(\i,3)--++(2,1.52);}
\foreach \j in {1,2}{\draw[tmBlue!60,line width=0.35pt] (0,\j)--(6,\j)--++(2,1.52);}
\foreach \k in {1,2,3}{\draw[tmBlue!60,line width=0.35pt] ($(0,3)+\k*(0.5,0.38)$)--++(6,0)--++(0,-3);}
\draw[geo] (Q)--(P)--(C)--(D)--cycle (D)--(A)--(B)--(C) (P)--(N)--(B);
\end{tikzpicture}
\end{minipage}
\end{probbox}
\probsol{
Mỗi lớp có $6\times4=24$ hình lập phương, có $3$ lớp.\par
Số hình lập phương: $24\times3=72$.\par
Thể tích: $72$ cm$^3$ (hoặc $6\times4\times3=72$ cm$^3$).
\dapso{$72$ hình lập phương; $72$ cm$^3$.}}
\gap
\begin{probbox}{7}{\relax}
Một hình hộp chữ nhật có thể tích $360$ cm$^3$, chiều dài $9$ cm, chiều rộng $8$ cm. Tính chiều cao của hình hộp.
\end{probbox}
\probsol{
Diện tích đáy: $9\times8=72$ cm$^2$.\par
Chiều cao: $360:72=5$ cm.
\dapso{$5$ cm.}}
\gap
\begin{probbox}{8}{\relax}
Một hình hộp chữ nhật có thể tích $1{,}8$ m$^3$, diện tích đáy $1{,}2$ m$^2$. Tính chiều cao của hình hộp.
\end{probbox}
\probsol{
$c=V:S_{\text{đáy}}=1{,}8:1{,}2=1{,}5$ m.
\dapso{$1{,}5$ m.}}
\gap
\begin{probbox}{9}{\relax}
Một hình lập phương có thể tích $512$ cm$^3$. Tính độ dài cạnh của hình lập phương.
\end{probbox}
\probsol{
Ta tìm số nhân với chính nó ba lần được $512$.\par
$7\times7\times7=343$, $8\times8\times8=512$. Cạnh dài $8$ cm.
\dapso{$8$ cm.}}
\gap
\begin{probbox}{10}{\relax}
Một bể nước dạng hình hộp chữ nhật dài $2$ m, rộng $1{,}5$ m, cao $1{,}2$ m. Khi đầy, bể chứa được bao nhiêu lít nước?
\end{probbox}
\probsol{
$V=2\times1{,}5\times1{,}2=3{,}6$ m$^3=3\,600$ dm$^3=3\,600$ l.
\dapso{$3\,600$ l.}}
\gap
\begin{probbox}{11}{\relax}
Một bể cá dài $60$ cm, rộng $30$ cm, cao $40$ cm. Mực nước bằng $\dfrac34$ chiều cao của bể. Trong bể có bao nhiêu lít nước?
\end{probbox}
\probsol{
Chiều cao mực nước: $40\times\dfrac34=30$ cm.\par
Thể tích nước: $60\times30\times30=54\,000$ cm$^3=54$ dm$^3=54$ l.
\dapso{$54$ l.}}
\gap
\begin{probbox}{12}{\relax}
Một bể nước dài $1{,}2$ m, rộng $0{,}9$ m, cao $1$ m. Lượng nước trong bể chiếm $60\%$ thể tích bể. Trong bể có bao nhiêu lít nước?
\end{probbox}
\probsol{
Thể tích bể: $1{,}2\times0{,}9\times1=1{,}08$ m$^3=1\,080$ l.\par
Lượng nước: $1\,080\times60:100=648$ l.
\dapso{$648$ l.}}
\gap
\begin{probbox}{13}{\relax}
Một bể có đáy dài $40$ cm, rộng $25$ cm. Thả vào bể một hòn đá ngập hoàn toàn trong nước thì mực nước dâng lên $3$ cm. Tính thể tích hòn đá theo đề-xi-mét khối.
\end{probbox}
\probsol{
Thể tích hòn đá: $40\times25\times3=3\,000$ cm$^3=3$ dm$^3$.
\dapso{$3$ dm$^3$.}}
\gap
\begin{probbox}{14}{\relax}
Một bể kính hình lập phương cạnh $5$ dm, mực nước trong bể cao $3$ dm. Sau khi thả vào bể một hòn đá (ngập hoàn toàn trong nước), mực nước cao $3{,}4$ dm. Tính thể tích hòn đá.
\end{probbox}
\probsol{
Phần nước dâng lên: $3{,}4-3=0{,}4$ dm.\par
Thể tích hòn đá: $5\times5\times0{,}4=10$ dm$^3$.
\dapso{$10$ dm$^3$.}}
\gap
\begin{probbox}{15}{\relax}
Một bể có đáy dài $50$ cm, rộng $30$ cm, mực nước cao $20$ cm. Thả vào bể một hòn đá có thể tích $1{,}8$ dm$^3$, ngập hoàn toàn trong nước. Hỏi lúc này mực nước cao bao nhiêu?
\end{probbox}
\probsol{
$1{,}8$ dm$^3=1\,800$ cm$^3$.\quad Diện tích đáy: $50\times30=1\,500$ cm$^2$.\par
Nước dâng lên: $1\,800:1\,500=1{,}2$ cm.\par
Mực nước lúc này: $20+1{,}2=21{,}2$ cm.
\dapso{$21{,}2$ cm.}}
\gap
\begin{probbox}{16}{\relax}
Một khối kim loại hình lập phương có cạnh $\dfrac15$ m. Mỗi xăng-ti-mét khối kim loại đó nặng $7{,}8$ g. Hỏi khối kim loại nặng bao nhiêu ki-lô-gam?
\end{probbox}
\probsol{
$\dfrac15$ m $=20$ cm.\quad Thể tích: $20\times20\times20=8\,000$ cm$^3$.\par
Khối lượng: $7{,}8\times8\,000=62\,400$ g $=62{,}4$ kg.
\dapso{$62{,}4$ kg.}}
\gap
\begin{probbox}{17}{\relax}
Một hình hộp chữ nhật dài $10$ cm, rộng $8$ cm, cao $6$ cm. Một hình lập phương có cạnh bằng trung bình cộng ba kích thước của hình hộp chữ nhật. Hình nào có thể tích lớn hơn và lớn hơn bao nhiêu?
\end{probbox}
\probsol{
Hình hộp chữ nhật: $10\times8\times6=480$ cm$^3$.\par
Cạnh hình lập phương: $(10+8+6):3=8$ cm.\quad Hình lập phương: $8\times8\times8=512$ cm$^3$.\par
$512-480=32$ cm$^3$.
\dapso{Hình lập phương lớn hơn $32$ cm$^3$.}}
\gap
\begin{probbox}{18}{\relax}
Hình lập phương B có cạnh gấp $2$ lần cạnh hình lập phương A. Hỏi thể tích hình lập phương B gấp bao nhiêu lần thể tích hình lập phương A?
\end{probbox}
\probsol{
Gọi cạnh hình A là $a$ thì cạnh hình B là $a\times2$.\par
$V_B=(a\times2)\times(a\times2)\times(a\times2)=a\times a\times a\times8=V_A\times8$.
\dapso{$8$ lần.}}
\gap
\begin{probbox}{19}{\relax}
Một bể không có nước dài $1{,}6$ m, rộng $1$ m, cao $1{,}2$ m. Một vòi chảy vào bể mỗi phút $40$ l nước. Hỏi sau bao lâu thì bể đầy?
\end{probbox}
\probsol{
Thể tích bể: $1{,}6\times1\times1{,}2=1{,}92$ m$^3=1\,920$ l.\par
Thời gian: $1\,920:40=48$ phút.
\dapso{$48$ phút.}}
\gap
\begin{probbox}{20}{\relax}
Một bể đầy nước dài $3$ m, rộng $2$ m, cao $1{,}5$ m. Mỗi ngày gia đình dùng hết $1\,800$ l nước. Hỏi lượng nước đó đủ dùng trong bao nhiêu ngày?
\end{probbox}
\probsol{
Lượng nước: $3\times2\times1{,}5=9$ m$^3=9\,000$ l.\par
Số ngày: $9\,000:1\,800=5$ ngày.
\dapso{$5$ ngày.}}
\gap
\begin{probbox}{21}{\relax}
Mỗi viên gạch dài $20$ cm, rộng $10$ cm, cao $5$ cm. Người ta xếp gạch khít nhau thành một khối dài $2$ m, rộng $1$ m, cao $0{,}5$ m. Hỏi khối đó có bao nhiêu viên gạch?
\end{probbox}
\probsol{
Một viên gạch: $20\times10\times5=1\,000$ cm$^3=1$ dm$^3$.\par
Khối gạch: $2\times1\times0{,}5=1$ m$^3=1\,000$ dm$^3$.\par
Số viên gạch: $1\,000:1=1\,000$ viên.
\dapso{$1\,000$ viên gạch.}}
\gap
\begin{probbox}{22}{\relax}
Một thùng các-tông dài $20$ cm, rộng $12$ cm, cao $8$ cm. Người ta xếp kín thùng bằng các hộp phấn, mỗi hộp có kích thước $4$ cm $\times\,3$ cm $\times\,2$ cm. Hỏi xếp được bao nhiêu hộp phấn?
\end{probbox}
\probsol{
Theo chiều dài: $20:4=5$; theo chiều rộng: $12:3=4$; theo chiều cao: $8:2=4$.\par
Số hộp: $5\times4\times4=80$.\par
\emph{Thử lại:} $20\times12\times8=1\,920$ cm$^3$, $4\times3\times2=24$ cm$^3$, $1\,920:24=80$.
\dapso{$80$ hộp phấn.}}
\gap
\begin{probbox}{23}{\relax}
Một hình hộp chữ nhật có diện tích xung quanh $480$ cm$^2$, chiều cao $8$ cm; chiều dài hơn chiều rộng $6$ cm. Tính thể tích của hình hộp.
\end{probbox}
\probsol{
Chu vi đáy: $480:8=60$ cm.\quad Nửa chu vi: $30$ cm.\par
Chiều dài: $(30+6):2=18$ cm.\quad Chiều rộng: $12$ cm.\par
$V=18\times12\times8=1\,728$ cm$^3$.
\dapso{$1\,728$ cm$^3$.}}
\gap
\begin{probbox}{24}{\relax}
Một bể nước có chu vi đáy $26$ dm, chiều cao $4$ dm; chiều dài hơn chiều rộng $3$ dm. Bể chứa được bao nhiêu lít nước?
\end{probbox}
\probsol{
Nửa chu vi: $26:2=13$ dm.\quad Chiều dài: $(13+3):2=8$ dm.\quad Chiều rộng: $5$ dm.\par
$V=8\times5\times4=160$ dm$^3=160$ l.
\dapso{$160$ l.}}
\gap
\begin{probbox}{25}{\relax}
Một bể kính hình lập phương cạnh $8$ dm, mực nước cao $5$ dm. Sau khi thả vào bể một cây san hô (ngập hoàn toàn trong nước), mực nước cao $5{,}5$ dm. Tính thể tích cây san hô.
\end{probbox}
\probsol{
Phần nước dâng lên: $5{,}5-5=0{,}5$ dm.\par
Thể tích cây san hô: $8\times8\times0{,}5=32$ dm$^3$.
\dapso{$32$ dm$^3$.}}
\gap
\begin{probbox}{26}{\relax}
Một bể nước dài $2{,}5$ m, rộng $1{,}6$ m, cao $1{,}2$ m. Mực nước trong bể cao $0{,}8$ m.\par
\pa{a} Trong bể có bao nhiêu lít nước?\par
\pa{b} Cần đổ thêm bao nhiêu lít nước nữa thì bể đầy?
\end{probbox}
\probsol{
\pa{a} $2{,}5\times1{,}6\times0{,}8=3{,}2$ m$^3=3\,200$ l.\par
\pa{b} Phần bể còn trống cao $1{,}2-0{,}8=0{,}4$ m; $2{,}5\times1{,}6\times0{,}4=1{,}6$ m$^3=1\,600$ l.
\dapso{\pa{a} $3\,200$ l; \pa{b} $1\,600$ l.}}
\gap
\begin{probbox}{27}{\relax}
Một bể nước dài $2$ m, rộng $1{,}5$ m, cao $1{,}4$ m. Lượng nước trong bể hiện chiếm $30\%$ thể tích bể. Hỏi phải đổ thêm bao nhiêu lít nước nữa để lượng nước trong bể chiếm $80\%$ thể tích bể?
\end{probbox}
\probsol{
Thể tích bể: $2\times1{,}5\times1{,}4=4{,}2$ m$^3=4\,200$ l.\par
Lượng nước cần đổ thêm bằng $80\%-30\%=50\%$ thể tích bể.\par
$4\,200\times50:100=2\,100$ l.
\dapso{$2\,100$ l.}}
\gap
\begin{probbox}{28}{\relax}
Một bể cá có đáy dài $50$ cm, rộng $40$ cm. Thả vào bể bốn hòn đá giống nhau, ngập hoàn toàn trong nước, thì mực nước dâng lên $1{,}5$ cm. Tính thể tích một hòn đá.
\end{probbox}
\probsol{
Bốn hòn đá: $50\times40\times1{,}5=3\,000$ cm$^3$.\par
Một hòn đá: $3\,000:4=750$ cm$^3$.
\dapso{$750$ cm$^3$.}}
\gap
\begin{probbox}{29}{\relax}
Một bể đầy nước. Người ta thả vào bể ba khối sắt hình lập phương giống nhau thì thấy $24$ l nước tràn ra. Tìm độ dài cạnh một khối sắt.
\end{probbox}
\probsol{
Ba khối sắt: $24$ l $=24$ dm$^3$.\quad Một khối: $24:3=8$ dm$^3$.\par
$2\times2\times2=8$ nên cạnh khối sắt dài $2$ dm.
\dapso{$2$ dm.}}
\gap
\begin{probbox}{30}{\relax}
Một hộp sữa dạng hình hộp chữ nhật dài $6$ cm, rộng $4$ cm, cao $10$ cm, được rót đầy sữa. Hỏi $4{,}8$ l sữa rót được đầy bao nhiêu hộp như thế?
\end{probbox}
\probsol{
Một hộp: $6\times4\times10=240$ cm$^3=240$ ml.\par
$4{,}8$ l $=4\,800$ ml.\quad Số hộp: $4\,800:240=20$ hộp.
\dapso{$20$ hộp.}}
\dividerbar{Thử thách -- Nâng cao}
\begin{probbox}{31}{\stars{1}}
Một bể cá dài $60$ cm, rộng $40$ cm, cao $45$ cm. Lúc đầu mực nước bằng $\dfrac35$ chiều cao của bể. Sau khi thả vào bể một hòn đá (ngập hoàn toàn trong nước), mực nước bằng $\dfrac23$ chiều cao của bể. Tính thể tích hòn đá.
\end{probbox}
\probsol{
Mực nước lúc đầu: $45\times\dfrac35=27$ cm.\quad Mực nước lúc sau: $45\times\dfrac23=30$ cm.\par
Nước dâng lên: $30-27=3$ cm.\par
Thể tích hòn đá: $60\times40\times3=7\,200$ cm$^3$.
\dapso{$7\,200$ cm$^3$.}}
\gap
\begin{probbox}{32}{\stars{1}}
Một bể cá bằng kính không có nắp dài $80$ cm, rộng $50$ cm, cao $50$ cm; mực nước cao $30$ cm.\par
\pa{a} Tính diện tích kính dùng để làm bể.\par
\pa{b} Thả vào bể một hòn đá có thể tích $16$ dm$^3$, ngập hoàn toàn trong nước. Hỏi lúc này mực nước cao bao nhiêu?
\end{probbox}
\probsol{
\pa{a} Xung quanh: $(80+50)\times2\times50=13\,000$ cm$^2$; đáy: $80\times50=4\,000$ cm$^2$.\par
Diện tích kính: $13\,000+4\,000=17\,000$ cm$^2$ $=170$ dm$^2$.\par
\pa{b} $16$ dm$^3=16\,000$ cm$^3$.\quad Nước dâng: $16\,000:4\,000=4$ cm.\quad Mực nước: $30+4=34$ cm.
\dapso{\pa{a} $170$ dm$^2$; \pa{b} $34$ cm.}}
\gap
\begin{probbox}{33}{\stars{2}}
Một bể dài $60$ cm, rộng $40$ cm, cao $50$ cm; mực nước cao $45$ cm. Người ta thả vào bể một hòn đá có thể tích $15$ dm$^3$. Hỏi có bao nhiêu lít nước tràn ra ngoài?
\end{probbox}
\probsol{
Phần bể còn trống phía trên mặt nước: $60\times40\times(50-45)=12\,000$ cm$^3=12$ dm$^3$.\par
Hòn đá ($15$ dm$^3$) lớn hơn phần còn trống nên nước tràn ra.\par
Nước tràn ra: $15-12=3$ dm$^3=3$ l.
\dapso{$3$ l.}}
\gap
\begin{probbox}{34}{\stars{1}}
Bể A dài $1{,}5$ m, rộng $1$ m, cao $0{,}8$ m, đang đầy nước. Người ta tháo hết nước ở bể A sang bể B (chưa có nước) có đáy dài $2$ m, rộng $1{,}2$ m. Hỏi mực nước ở bể B cao bao nhiêu?
\end{probbox}
\probsol{
Lượng nước: $1{,}5\times1\times0{,}8=1{,}2$ m$^3$.\par
Diện tích đáy bể B: $2\times1{,}2=2{,}4$ m$^2$.\par
Mực nước: $1{,}2:2{,}4=0{,}5$ m.
\dapso{$0{,}5$ m.}}
\gap
\begin{probbox}{35}{\stars{1}}
Diện tích toàn phần của một hình lập phương lớn hơn diện tích xung quanh của nó $72$ cm$^2$. Tính thể tích hình lập phương đó.
\end{probbox}
\probsol{
Phần hơn là diện tích hai mặt đáy: một mặt là $72:2=36$ cm$^2=6\times6$.\par
Cạnh $6$ cm; $V=6\times6\times6=216$ cm$^3$.
\dapso{$216$ cm$^3$.}}
\gap
\begin{probbox}{36}{\stars{1}}
Một hình hộp chữ nhật có thể tích $120$ cm$^3$. Nếu gấp chiều dài lên $2$ lần, gấp chiều rộng lên $3$ lần và giảm chiều cao đi $2$ lần thì thể tích mới là bao nhiêu?
\end{probbox}
\probsol{
Thể tích được nhân với $2\times3=6$ rồi chia cho $2$, tức là gấp lên $3$ lần.\par
Thể tích mới: $120\times3=360$ cm$^3$.
\dapso{$360$ cm$^3$.}}
\gap
\begin{probbox}{37}{\stars{2}}
Cạnh của một hình lập phương dài $10$ cm được tăng thêm $20\%$. Hỏi thể tích tăng thêm bao nhiêu xăng-ti-mét khối và bao nhiêu phần trăm?
\end{probbox}
\probsol{
Cạnh mới: $10+10\times20:100=12$ cm.\par
Thể tích: $10\times10\times10=1\,000$ cm$^3$ và $12\times12\times12=1\,728$ cm$^3$.\par
Tăng thêm: $1\,728-1\,000=728$ cm$^3$, tức là $728:1\,000=0{,}728=72{,}8\%$.
\dapso{$728$ cm$^3$; $72{,}8\%$.}}
\gap
\begin{probbox}{38}{\stars{2}}
Ba mặt khác nhau của một hình hộp chữ nhật có diện tích $24$ cm$^2$, $30$ cm$^2$ và $20$ cm$^2$; các cạnh có số đo là số tự nhiên (xăng-ti-mét). Tính thể tích hình hộp.
\end{probbox}
\probsol{
Cạnh chung của hai mặt $24$ cm$^2$ và $30$ cm$^2$ chia hết cả $24$ và $30$: có thể là $1$; $2$; $3$ hoặc $6$.\par
Thử $6$: hai cạnh còn lại là $24:6=4$ và $30:6=5$, và $4\times5=20$. \checkmark\par
(Với $1$; $2$; $3$ thì mặt thứ ba là $720$; $180$; $80$, không bằng $20$.)\par
$V=6\times4\times5=120$ cm$^3$.
\dapso{$120$ cm$^3$.}}
\gap
\begin{probbox}{39}{\stars{2}}
Một hình hộp chữ nhật có chiều dài gấp $3$ lần chiều cao, chiều rộng gấp $2$ lần chiều cao. Thể tích của hình hộp là $384$ cm$^3$. Tính các kích thước của hình hộp.
\end{probbox}
\probsol{
Chia hình hộp thành các hình lập phương có cạnh bằng chiều cao: được $3\times2\times1=6$ hình.\par
Một hình lập phương: $384:6=64$ cm$^3=4\times4\times4$, nên chiều cao là $4$ cm.\par
Chiều dài: $12$ cm; chiều rộng: $8$ cm.
\dapso{$12$ cm, $8$ cm, $4$ cm.}}
\gap
\begin{probbox}{40}{\stars{2}}
Một bể chưa có nước dài $1{,}5$ m, rộng $1{,}2$ m, cao $1$ m. Vòi A chảy vào bể mỗi phút $30$ l, vòi B tháo nước ra mỗi phút $12$ l. Mở cả hai vòi cùng lúc thì sau bao lâu bể đầy?
\end{probbox}
\probsol{
Thể tích bể: $1{,}5\times1{,}2\times1=1{,}8$ m$^3=1\,800$ l.\par
Mỗi phút lượng nước trong bể tăng: $30-12=18$ l.\par
Thời gian: $1\,800:18=100$ phút $=1$ giờ $40$ phút.
\dapso{$1$ giờ $40$ phút.}}
\gap
\begin{probbox}{41}{\stars{1}}
Một bể bơi dài $20$ m, rộng $10$ m. Một máy bơm mỗi giờ bơm vào bể $60$ m$^3$ nước. Hỏi sau bao lâu thì mực nước trong bể (lúc đầu chưa có nước) cao $1{,}5$ m?
\end{probbox}
\probsol{
Lượng nước cần bơm: $20\times10\times1{,}5=300$ m$^3$.\par
Thời gian: $300:60=5$ giờ.
\dapso{$5$ giờ.}}
\gap
\begin{probbox}{42}{\stars{2}}
Trong một trận mưa, lượng mưa đo được là $15$ mm trên một mái nhà bằng dài $8$ m, rộng $6$ m. Toàn bộ nước mưa chảy vào một bể chứa. Hỏi bể nhận thêm bao nhiêu lít nước?
\end{probbox}
\probsol{
$15$ mm $=0{,}015$ m.\par
Nước mưa trên mái: $8\times6\times0{,}015=0{,}72$ m$^3=720$ l.
\dapso{$720$ l.}}
\gap
\begin{probbox}{43}{\stars{1}}
\begin{minipage}[t]{0.6\linewidth}
Một khối gỗ hình lập phương cạnh $10$ cm. Người ta khoét một lỗ vuông $4$ cm $\times\,4$ cm xuyên thẳng từ mặt trên xuống mặt dưới. Tính thể tích phần gỗ còn lại.
\end{minipage}\hfill
\begin{minipage}[t]{0.36\linewidth}\vspace{0pt}\centering
\begin{tikzpicture}[scale=0.17]
\cuboidVT{10}{10}{10}{0}
\coordinate (h1) at ($(D)+(3,0)+(1.5,1.14)$);
\fill[white,draw=tmBlue!75!black,line width=0.6pt] (h1)--++(4,0)--++(2,1.52)--++(-4,0)--cycle;
\node[dim,below] at ($(Q)!0.5!(P)$){$10$ cm};
\end{tikzpicture}
\end{minipage}
\end{probbox}
\probsol{
Khối lập phương: $10\times10\times10=1\,000$ cm$^3$.\par
Lỗ khoét là hình hộp chữ nhật: $4\times4\times10=160$ cm$^3$.\par
Phần gỗ còn lại: $1\,000-160=840$ cm$^3$.
\dapso{$840$ cm$^3$.}}
\gap
\begin{probbox}{44}{\stars{2}}
Khi nước đóng băng, thể tích tăng thêm $\dfrac1{11}$. Một khay làm đá chứa lượng nước dài $22$ cm, rộng $10$ cm, sâu $3$ cm. Tính thể tích khối đá sau khi nước đóng băng.
\end{probbox}
\probsol{
Thể tích nước: $22\times10\times3=660$ cm$^3$.\par
Phần tăng thêm: $660\times\dfrac1{11}=60$ cm$^3$.\par
Thể tích khối đá: $660+60=720$ cm$^3$.
\dapso{$720$ cm$^3$.}}
\gap
\begin{probbox}{45}{\stars{1}}
Một bể nước có đáy dài $2$ m, rộng $1{,}2$ m. Buổi sáng mực nước cao $1$ m, đến chiều mực nước còn $0{,}65$ m. Hỏi trong ngày đã dùng bao nhiêu lít nước?
\end{probbox}
\probsol{
Mực nước giảm: $1-0{,}65=0{,}35$ m.\par
Lượng nước đã dùng: $2\times1{,}2\times0{,}35=0{,}84$ m$^3=840$ l.
\dapso{$840$ l.}}
\gap
\begin{probbox}{46}{\stars{2}}
Có thể xếp được nhiều nhất bao nhiêu hình lập phương cạnh $3$ cm vào một cái hộp dài $16$ cm, rộng $12$ cm, cao $9$ cm (các hình đặt ngay ngắn, sát nhau)? Phần thể tích còn trống của hộp là bao nhiêu?
\end{probbox}
\probsol{
Theo chiều dài: $16:3=5$ (dư $1$) nên xếp được $5$ hình; theo chiều rộng: $12:3=4$; theo chiều cao: $9:3=3$.\par
Số hình lập phương: $5\times4\times3=60$.\par
Hộp: $16\times12\times9=1\,728$ cm$^3$; các hình lập phương: $60\times27=1\,620$ cm$^3$.\par
Phần còn trống: $1\,728-1\,620=108$ cm$^3$.
\dapso{$60$ hình; $108$ cm$^3$.}}
\gap
\begin{probbox}{47}{\stars{1}}
Một hình hộp chữ nhật có thể tích $1{,}2$ m$^3$, diện tích đáy $0{,}8$ m$^2$, chiều dài $1$ m. Tính chiều cao và diện tích xung quanh của hình hộp.
\end{probbox}
\probsol{
Chiều cao: $1{,}2:0{,}8=1{,}5$ m.\quad Chiều rộng: $0{,}8:1=0{,}8$ m.\par
$S_{xq}=(1+0{,}8)\times2\times1{,}5=5{,}4$ m$^2$.
\dapso{$1{,}5$ m; $5{,}4$ m$^2$.}}
\gap
\begin{probbox}{48}{\stars{1}}
Người ta tháo hết nước từ một bể hình lập phương cạnh $6$ dm đang đầy nước sang một bể chưa có nước có đáy dài $9$ dm, rộng $8$ dm. Hỏi mực nước trong bể sau cao bao nhiêu?
\end{probbox}
\probsol{
Lượng nước: $6\times6\times6=216$ dm$^3$.\par
Diện tích đáy: $9\times8=72$ dm$^2$.\par
Mực nước: $216:72=3$ dm.
\dapso{$3$ dm.}}
\gap
\begin{probbox}{49}{\stars{2}}
\begin{minipage}[t]{0.6\linewidth}
Một cầu thang bê tông có $3$ bậc, rộng $1{,}5$ m. Mỗi bậc sâu $30$ cm, cao $15$ cm (xem hình nhìn từ bên cạnh). Tính thể tích bê tông.
\end{minipage}\hfill
\begin{minipage}[t]{0.36\linewidth}\vspace{0pt}\centering
\begin{tikzpicture}[scale=0.032]
\draw[geo,fill=tmBlue!10] (0,0)--(90,0)--(90,45)--(60,45)--(60,30)--(30,30)--(30,15)--(0,15)--cycle;
\draw[tmBlue!60,line width=0.35pt] (30,0)--(30,15) (60,0)--(60,30) (30,15)--(90,15) (60,30)--(90,30);
\node[dim,below] at (15,0){$30$ cm}; \node[dim,right] at (90,7.5){$15$ cm};
\end{tikzpicture}
\end{minipage}
\end{probbox}
\probsol{
Chia khối bậc thang thành ba lớp, mỗi lớp cao $15$ cm $=0{,}15$ m:\par
lớp dưới sâu $0{,}9$ m, lớp giữa sâu $0{,}6$ m, lớp trên sâu $0{,}3$ m.\par
Thể tích: $1{,}5\times0{,}15\times(0{,}9+0{,}6+0{,}3)=1{,}5\times0{,}15\times1{,}8=0{,}405$ m$^3$.
\dapso{$0{,}405$ m$^3$.}}
\gap
\begin{probbox}{50}{\stars{2}}
Một bể nước dài $1{,}5$ m, mực nước trong bể cao $60$ cm. Đổ thêm $240$ l nước thì mực nước cao $80$ cm. Tính chiều rộng của bể.
\end{probbox}
\probsol{
$240$ l $=240$ dm$^3$; nước dâng thêm $80-60=20$ cm $=2$ dm.\par
Diện tích đáy: $240:2=120$ dm$^2=1{,}2$ m$^2$.\par
Chiều rộng: $1{,}2:1{,}5=0{,}8$ m.
\dapso{$0{,}8$ m.}}
